\documentclass[10pt,a4paper]{article}

\usepackage[T1]{fontenc}
\usepackage[italian]{babel}
\usepackage{amsmath,amssymb,amsfonts}
\usepackage{geometry}
\usepackage{tabularx}
\usepackage{booktabs}
\usepackage{array}
\usepackage[table]{xcolor}
\usepackage{tikz}

\definecolor{lightgray}{gray}{0.85}

\newcommand{\grigioriga}{\arrayrulecolor{lightgray}\midrule\arrayrulecolor{black}}

\geometry{
    a4paper,
    left=15mm,
    right=15mm,
    top=20mm,
    bottom=20mm
}

\newcolumntype{M}{>{\centering\arraybackslash$\displaystyle}X<{$}}
\newcolumntype{L}{>{\raggedright\arraybackslash\bfseries}p{5cm}}

\title{\Huge\textbf{Formulario Fisica}}
\author{}
\date{}

\begin{document}

\maketitle

\section{Cinematica}

\subsection{Leggi orarie}

\begin{tabularx}{\textwidth}{L M}
    \toprule
    \large Descrizione & \multicolumn{1}{c}{\large Formula} \\ 
    
    \midrule
    Spostamento & x(t) = x_0 + v_0 t + \frac{1}{2} a t^2 \\[1.5ex]
    \grigioriga
    Velocità & v(t) = v_0 + a(t - t_0) \\[1.5ex]
    \grigioriga
    Velocità in Posizione & v^2(x) = v_0^2 + 2a(x - x_1) \\
    \grigioriga
    Media & v = \frac{v_1 - v_0}{v_1 - t_0} \ \text{(analogo per accelerazione)}\\[1.5ex]
    \grigioriga
    Velocità con accelerazione & v(t) = v_0 + \int_{t_0}^{t_1} a(t) \ dt \\[1.5ex]
    
    \bottomrule
\end{tabularx}

\subsection{Moto del proiettile}

\begin{tabularx}{\textwidth}{L M}
    \toprule
    \large Nome Formula & \multicolumn{1}{c}{\large Formula} \\ 
    
    \midrule
    Posizione x & x(t) = x_0 + (v_0 \cos\phi_0)t \\[1.5ex]
    \grigioriga
    Posizione y & y(t) = y_0 + (v_0 \sin\phi_0)t - \frac{1}{2}gt^2 \\[1.5ex]
    \grigioriga
    Modulo velocità & |\vec{v}| = \sqrt{v_x^2 + v_y^2} \\[1.5ex]
    \grigioriga
    Angolo iniziale & tan\phi_0 = \frac{v_y}{v_x} \\[1.5ex]
    \grigioriga
    Gittata & R = \frac{v_0^2}{g} \sin 2\phi_0 \\[1.5ex] 

    \bottomrule
\end{tabularx}

\subsection{Moto circolare}

\begin{tabularx}{\textwidth}{L M}
    \toprule
    \large Nome Formula & \multicolumn{1}{c}{\large Formula} \\
    
    \midrule
    Posizione x & x = R \cos \theta \\[1.5ex]
    \grigioriga
    Posizione y & y = R \sin \theta \\[1.5ex]
    \grigioriga
    Vettore posizione & \vec{r} = R \cos\theta \hat{i} + R \sin\theta \hat{j} \\[1.5ex]
    \grigioriga
    Velocità angolare & \omega = \frac{v_s}{R} \\[1.5ex]
    \grigioriga
    Angolo orario & \theta(t) = \theta_0 + \omega(t - t_0) \\[1.5ex]
    \grigioriga
    Accelerazione centripeta & a_c = \omega^2 R = \frac{v_s^2}{R} = v_s \omega \\[1.5ex]
    \grigioriga
    Accelerazione angolare & \alpha = \frac{d\omega}{dt} \\[1.5ex]
    \grigioriga
    Accelerazione tangenziale & a_t = \alpha R\\[1.5ex]
    \grigioriga
    Accelerazione totale & a = \sqrt{a_c^2 + a_t^2} = R\sqrt{\omega^4 + \alpha^2} \\[1.5ex]
    
    \bottomrule
\end{tabularx}

\section{Forze}

\begin{tabularx}{\textwidth}{L M}
    \toprule
    \large Nome Formula & \multicolumn{1}{c}{\large Formula} \\
    
    \midrule
    2° Legge di Newton & \vec{F} = m\vec{a} \\[1.5ex]
    \grigioriga
    Forza peso & \vec{P} = m\vec{g} \\[1.5ex]
    \grigioriga
    Forza vettoriale & \vec{F} = F\cos\theta + F\sin\theta \\[1.5ex]
    \grigioriga
    Forza di attrito statica & F_s \le \mu_s N \\[1.5ex]
    \grigioriga
    Forza di attrito dinamica & F_k = \mu_k N  \\[1.5ex]
    \grigioriga
    Forza elastica & F = -k (l - l_0) \ \text{con $k$ costante elastica} \\[1.5ex]
    \grigioriga
    Oscillatore armonico & x(t) = A \cos(\omega t + \phi) \ \text{con $\omega = \sqrt{\frac{k}{m}}$, $A$ Ampiezza, $\omega$ Pulsazione, $\phi$ Fase} \\[1.5ex] 
    \grigioriga
    Periodo oscillatore armonico & T = \frac{2\pi}{\omega} = 2\pi \sqrt{\frac{m}{k}} \\[1.5ex]
    \grigioriga
    Frequenza oscillatore armonico & \nu = \frac{1}{T} = \frac{\omega}{2\pi} = \frac{1}{2\pi} \sqrt{\frac{k}{m}} \\[1.5ex]
    \grigioriga
    Velocità max moto armonico & v_{\text{max}} = \omega A \\[1.5ex]
    \grigioriga
    Accelerazione max moto armonico & a_\text{max} = \omega^2 A\\[1.5ex]
    \grigioriga
    Pendolo semplice & \sum F_x = -m g \sin \theta = m a \implies a = -g \sin \theta \\[1.5ex]
    Piccole oscillazioni pendolo & 2 \pi \sqrt{\frac{l}{g}}\\[1.5ex]
    Trazione & T = -\vec{P} + \vec{F}_\text{risalita} = mg + m \frac{v^2}{l} \\[1.5ex]
    
    \bottomrule
\end{tabularx}

\subsection{Lavoro ed Energia}{}

\begin{tabularx}{\textwidth}{L M}
    \toprule
    \large Nome Formula & \multicolumn{1}{c}{\large Formula} \\
    
    \midrule
    Lavoro & L = 
        \int_{\vec{r}_i}^{\vec{r}_f} \vec{F} \ d\vec{r} = 
        \frac{1}{2} m v_f^2 - \frac{1}{2} m v_i^2 = 
        F \frac{\Delta x}{\sin(\Theta)} \cos(\frac{\pi}{2} - \Theta) \\[1.5ex]
    \grigioriga
    Energia cinetica & K = \frac{1}{2} m v^2 \\[1.5ex]
    \grigioriga
    Energia potenziale & \text{d}U = -F \text{d}x\\[1.5ex]
    \grigioriga
    Energia meccanica & E = U + K \\[1.5ex]
    \grigioriga
    Energia potenziale peso & U(y) = mgy \\[1.5ex]
    \grigioriga
    Energia potenziale elastica & U(x) = \frac{1}{2} k x^2 \\[1.5ex]
    \grigioriga
    Energia meccanica oscillatore armonico & E = k + U = \frac{1}{2} m v^2 + \frac{1}{2} k x^2 \\[3.5ex]
    (in equilibrio) & U(0) = 0 \text{, } K = E\\[1.5ex]
    (al massimo) & U(A) = E \text{, } K = 0 \text{, con A ampiezza}\\[1.5ex]
    \grigioriga
    Potenza & P = \frac{dL}{dt} \\[1.5ex]
    Potenza su punto materiale & P = \vec{F} \dot \vec{v} \\[1.5ex]

    \bottomrule
\end{tabularx}

\section{Elettrostatica}
\begin{tabularx}{\textwidth}{L M}
    \toprule
    \large Nome Formula & \multicolumn{1}{c}{\large Formula} \\
    
    \midrule
    Carica oggetto & q = (N_p - N_e)e \\[1.5ex]
    Costante di proporzionalità & K = \frac{1}{4 \pi \epsilon_0 [\epsilon_r]} \text{ con $\epsilon_r$ solo non nel vuoto} \\[1.5ex]
    Coulomb, cariche puntiformi & \vec{F_{ab}} = \frac{1}{4 \pi \epsilon_0} \frac{q_a q_b}{r^2} \hat{r}_{ab} \text{ , } \vec{E} = \frac{1}{4 \pi \epsilon_0} \frac{q}{r^2} \hat{r} \\[1.5ex]
    Accelerazione da campo & a = \frac{q\vec{E}}{m} \\[1.5ex]
    Gauss & \Phi(\vec{E}) = \frac{Q}{\epsilon_0} \\[1.5ex]
    Potenziale per carica & V = \frac{1}{4 \pi \epsilon_0} \frac{q}{r} \\[1.5ex]
    Energia potenziale elettrica & U = V q \\[1.5ex]
    Capacità & C = \frac{Q}{V} \\[1.5ex]
    Condensatore piano & C = \frac{[\epsilon_r] \epsilon_0 A}{d} \\[1.5ex]
    Condensatori in serie & \frac{1}{C} = \sum_i \frac{1}{  C_i} \\[1.5ex]    
    Condensatori in parallelo & C = \sum_i C_i \\[1.5ex]
    Energia potenziale condensatore & U = \frac{Q^2}{2C} = \frac{1}{2} C V^2 = \frac{QV}{2} \\[1.5ex]   
    Densità di energia & u = \frac{U}{A d} = \frac{1}{2} \epsilon_0 E^2 \\[1.5ex]
    Intensità di corrente & I = \frac{dQ}{dt} = n S v_d |q| \text{, con $n$ densità portatori, $S$ sezione filo, $v_d$ velocità deriva} \\[1.5ex]
    Densità volumica & n = \frac{dN}{dV} \text{, con $N$ numero di particelle e $V$ volume} \\[1.5ex]
    Densità corrente & J = \frac{I}{S} = n q v_d \\[1.5ex]
    Legge di Ohm & V = R I \\[1.5ex]
    Resistività materiale & R = \rho \frac{l}{S} \\[1.5ex]
    Ohm con densità corrente & \vec{E} = \rho \vec{J} \\[1.5ex]
    Drude accelerazione con campo & \vec{a} = - \frac{e \vec{E}}{m} = \frac{\vec{F}}{m} \\[1.5ex]
    Drude velocità deriva & \vec{v_d} = -\frac{e E}{m} \tau \text{ con $\tau$ tempo medio tra urti} \\[1.5ex]
    Drude resistività & \rho = \frac{m}{n e^2 \tau} \\[1.5ex]
    Resistenze in serie & R_{tot} = R_1 + R_2 \\[1.5ex]
    Resistenze in parallelo & \frac{1}{R_{tot}} = \frac{1}{R_1} + \frac{1}{R_2} \\[1.5ex]
    Legge di Joule & P = V I = R I^2 = \frac{V^2}{R} \\[1.5ex]
    Energia dissipata in resistenza & \Delta U = -V \Delta Q \\[1.5ex]
    1° Legge Kirchoff & I_{in_1} + I_{in_2} = I_{out_1} + I_{out_2} \\[1.5ex]
    2° Legge Kirchoff & V - R_1 I - R_2 I - R_3 I = 0 \\[1.5ex]
    Carica condensatore & q(t) = C \Delta V ( 1 - e ^ {- \frac{t}{RC}}) \\[1.5ex]
    \bottomrule
\end{tabularx}

\section{Magnetismo}
\begin{tabularx}{\textwidth}{L M}
    \toprule
    \large Nome Formula & \multicolumn{1}{c}{\large Formula} \\
    
    \midrule
    Campo generato da un filo & B \propto \frac{i}{r} \text{, con $i$ corrente nel filo, $r$ distanza dal filo} \\[1.5ex]
    Legge di Lorentz & \vec{F} = q \vec{V} \cdot \vec{B} \\[1.5ex]
    Lorents con campo E & \vec{F} = q \vec{E} + q (\vec{v} \cdot \vec{B}) \\[1.5ex]
    Modulo Lorenz & F = |q| v B \sin \theta \text{, con $\theta$ angolo tra velocità e campo} \\[1.5ex]
    Moto campo magnetico uniforme & r = \frac{m v}{q B} \text{, } \omega_C = \frac{q B}{m} \text{ (detta frequenza di ciclotrone), } \nu = \frac{\omega_C}{2 \pi} \\[1.5ex]
    Velocità in ciclotrone & v = \frac{E}{B} \\[1.5ex]
    Thomson & \frac{e}{m} = \frac{y E}{B^2 (LD + \frac{L^2}{2})} \\[1.5ex]
    Campo di Hall & E_H = \frac{J B}{n q} = v_d B\\[1.5ex]
    Coefficiente di Hall & R_H = \frac{E_H}{JB} = \frac{1}{n q} \\[1.5ex]
    Filo rettilineo & F = I l \times B\\[1.5ex]
    Momento meccanico torcente & \tau = I S \times B \\[1.5ex]
    Teorema di Ampére & \oint B \ dr = \mu_0 \sum I_{conc}  \\[1.5ex]
    Filo dritto infinito & B = \frac{\mu_0 I}{2 \pi R} \\[1.5ex]
    Interno del filo & B = \frac{\mu_0 I R}{2 \pi a^2} \\[1.5ex]
    Solenoide & B = \mu_0 n I \text{ all'interno, all'esterno è nullo} \\[1.5ex]
    Toroide (ciambella) & B = \frac{\mu_0 I N}{2 \pi R} \text{, con $N$ numero spire} \\[1.5ex]
    Forza tra fili paralleli & F = \frac{\mu_0 I_1 I_2 l}{2 \pi R} \\[1.5ex]
    1° Laplace  & dB = \frac{\mu_0 I d l \times \hat{r}}{4 \pi r^2} \\[1.5ex]
    1° Laplace con 1 spira & B = \frac{\mu_0 I}{2 R} \\[1.5ex]
    Ferromagnetismo & B = \mu_r \mu_0 n I \text{, con $\mu_r$ la permeabilità del materiale} \\[1.5ex]
    \bottomrule
\end{tabularx}

\section{Magnetismo 2}
\begin{tabularx}{\textwidth}{L M}
    \toprule
    \large Nome Formula & \multicolumn{1}{c}{\large Formula} \\
    
    \midrule
    Flusso & \Phi(B) = B S \cos \theta \\[1.5ex]
    Faraday (Forza Elettromotrice Indotta) & \mathcal{E}_{ind} = - \frac{d}{dt} (\Phi(B)) \text{, il flusso cambia nel tempo per $B$, $S$ o $\theta$}\\[1.5ex]
    Corrente da F.E.M. con Ohm & I = \frac{| \mathcal{E}_{ind} |}{R} \\[1.5ex]
    Verso corrente & \text{Se flusso fuori \textbf{aumenta} viene \textbf{respinto}, se fuori \textbf{diminuisce} viene \textbf{aiutato}} \\[1.5ex]
    Induzione di moto & \mathcal{E}_{ind} = B L_\text{filo} v_\text{filo} \\[1.5ex]
    Velocità limite & v_\text{max} = \frac{m g R}{B^2 L_\text{filo}^2} \\[2.5ex]
    F.E.M. Alternatore & \mathcal{E}_\text{ind} = N B S \omega \sin(\omega t) \text{, con $S$ area, N spire, ruota a velocità $\omega$} \\[1.5ex]
     & \mathcal{E}_\text{max} = N B S \omega \text{ (ovvero l'Ampiezza) }\\[1.5ex]
    F.E.M. Trasformatore & \frac{V_P}{V_S} = \frac{N_P}{N_S} \text{, } V_P I_P = V_S I_S \\[1.5ex]
    Corrente di Spostamento (Ampère-Maxwell) & I_s = \epsilon_0 \frac{d}{dt} \Phi (E) \\[1.5ex]
    
    \bottomrule
\end{tabularx}

\section{Trigonometria}


\definecolor{lato_colore}{RGB}{30, 41, 59}   
\definecolor{angolo_alpha}{RGB}{220, 38, 38} 
\definecolor{angolo_beta}{RGB}{29, 78, 216}  
\definecolor{angolo_retto}{RGB}{21, 128, 61} 


\begin{figure}[htbp]
    \begin{minipage}[c]{0.55\textwidth}
        \begin{tikzpicture}[scale=1.5, thick, >=stealth]
            \coordinate (C) at (0,0);
            \coordinate (B) at (5,0);
            \coordinate (A) at (0,3.5);

            \draw[fill=angolo_retto!20, draw=angolo_retto, thick] (0,0) -- (0.4,0) -- (0.4,0.4) -- (0,0.4) -- cycle;
            
            \begin{scope}
                \clip (A) -- (C) -- (B) -- cycle;
                \draw[fill=angolo_alpha!20, draw=angolo_alpha, thick] (A) circle (0.7cm);
            \end{scope}
            
            \begin{scope}
                \clip (B) -- (A) -- (C) -- cycle;
                \draw[fill=angolo_beta!20, draw=angolo_beta, thick] (B) circle (0.8cm);
            \end{scope}

            \draw[color=lato_colore, line width=1.2pt] (C) -- (B) node[midway, below=4pt, font=\bfseries] {$a$}
                -- (A) node[midway, above right=2pt, font=\bfseries] {$c$}
                -- cycle node[midway, left=4pt, font=\bfseries] {$b$};
            \node[below left, color=lato_colore] at (C) {$C$};
            \node[below right, color=lato_colore] at (B) {$B$};
            \node[above left, color=lato_colore] at (A) {$A$};

            \node[color=angolo_alpha] at (0.25, 2.7) {$\alpha$};
            \node[color=angolo_beta] at (4.1, 0.25) {$\beta$};
            \node[color=angolo_retto] at (0.5, 0.5) {$\gamma = 90^\circ$};
        \end{tikzpicture} 
    \end{minipage}
    \begin{minipage}[c]{0.4\textwidth}

        \begin{tabularx}{\textwidth}{M}
        a = c \cdot {\color{angolo_beta}\cos \beta} \\
        a = c \cdot {\color{angolo_alpha}\sin \alpha} \\[1.5ex]
        
        b = c \cdot {\color{angolo_alpha}\cos \alpha} \\
        b = c \cdot {\color{angolo_beta}\sin \beta} \\[1.5ex]
        
        c = \frac{b}{{\color{angolo_beta}\sin \beta}} = \frac{a}{{\color{angolo_alpha}\sin \alpha}} = \frac{a}{{\color{angolo_beta}\cos \beta}} = \frac{b}{{\color{angolo_alpha}\cos \alpha}}
    \end{tabularx}

    \end{minipage}
\end{figure}

\section{Costanti}
\begin{tabularx}{\textwidth}{L M}

    \toprule
    \large Nome Formula & \multicolumn{1}{c}{\large Formula} \\

    \midrule
    Carica elementare & e = 1.60218 \cdot 10^{-19} C \\[1.5ex]
    Costante dielettrica vuoto & \epsilon_0 = 8.854 \cdot 10^{-12} \frac{C^2}{Nm^2} \\[1.5ex]
    Avogadro & N_a = 6.022 \cdot 10^{23} \\[1.5ex]
    Massa elettrone & m_e = 9.11 \cdot 10^{-31} kg \\[1.5ex]
    Massa protone & m_p = 1.67 \cdot 10^{-27} kg \\[1.5ex]
    Permeabilità magnetica vuoto & \mu_0 = 4 \pi \cdot 10^{-7} \ \frac{T_m}{A} \\[1.5ex]
 
    \bottomrule

\end{tabularx}

\end{document}